\section{Conclusion and Future Work}
\label{sec:conclusion}

We have presented a structured research roadmap towards a rigorous proof of the Yang-Mills Mass Gap Conjecture. By explicitly distinguishing between established rigorous results and the remaining proof obligations, we provide a clear path forward for the mathematical physics community.

\subsection{Summary of Rigorous Foundations}

This work solidifies the following foundations:
\begin{enumerate}
    \item \textbf{Lattice Definition:} The theory is rigorously defined on the lattice, satisfying Osterwalder-Schrader reflection positivity.
    \item \textbf{Strong Coupling Gap:} The mass gap is proven to exist for $\beta < \beta_0$ using convergent cluster expansions.
    \item \textbf{Finite Volume Gap:} The gap is strictly positive in any finite volume.
\end{enumerate}

\subsection{The Path Forward: Resolving the Blockers}

To complete the proof, the following ``work packages'' must be executed:

\begin{enumerate}
    \item \textbf{Intermediate Coupling (The ``Gap''):} We have proposed the \textbf{Adjoint Interpolation} strategy (Strategy B) as a promising route. The goal is to rigorously prove that the path from $\mathcal{N}=1$ SYM (which is gapped) to pure Yang-Mills (infinite mass limit) encounters no phase transitions. This relies on establishing the analyticity of the free energy and the protection of the gap by center symmetry.
    
    \item \textbf{Continuum Limit:} We have outlined the application of Balaban's renormalization group methods. The remaining task is to verify the technical conditions required for the specific gauge group and action, ensuring the axioms hold in the limit $a \to 0$.
    
    \item \textbf{Physical Scaling:} We must rigorously prove the scaling inequality $\Delta \ge c \sqrt{\sigma}$ in the continuum limit to ensure the physical gap does not vanish.
\end{enumerate}

\subsection{Final Assessment}

While the Yang-Mills Mass Gap problem remains open, this roadmap identifies the precise mathematical obstacles and proposes concrete, non-perturbative strategies to overcome them. We believe that the ``Adjoint Interpolation'' method offers the most viable route to a solution, transforming the problem of ``proving a gap from scratch'' into a problem of ``deformation and analyticity'' from a known supersymmetric limit.
